%!LW recipe=latexmk-xelatex
\documentclass[UTF8,zihao=-4]{ctexart}

% 页面布局，四周保留 7 mm 页边距
\usepackage[paperwidth=250mm,paperheight=245mm,margin=7mm]{geometry}
\usepackage{array,multirow,graphicx}

% 字体：正文-仿宋_GB2312;数字字母-Times New Roman;标题-方正小标宋简体
\usepackage{fontspec}
\setmainfont{Times New Roman}
\setCJKmainfont[AutoFakeBold=2.2]{FangSong_GB2312}
\newCJKfontfamily\titlefont{FZXiaoBiaoSong-B05S}

% 页面与表格
\pagestyle{empty}
\setlength{\parindent}{0pt}
\setlength{\tabcolsep}{2.5pt}
\setlength{\arrayrulewidth}{0.45pt}
\renewcommand{\arraystretch}{1.18}

% 列格式：水平、垂直居中
\newcolumntype{C}[1]{>{\centering\arraybackslash}m{#1}}
% 合并三列并居中
\newcommand{\cell}[1]{\multicolumn{3}{c|}{#1}}

\begin{document}
\begin{center}

% 标题居中，日期靠右
\makebox[\linewidth][c]{%
  {\zihao{2}\titlefont 北京航空航天大学 系号--学院/书院对照表}}%
\makebox[0pt][r]{{\zihao{4}(2026.7)}}\\[2mm]

% 表格拆为八个基础列；普通学院行通过 multicolumn 合并，分组行直接使用基础列，避免嵌套表格导致横竖线缺失
\begin{tabular}{
|C{11mm}|C{13mm}|C{15mm}|C{68mm}|
 C{12mm}|C{15mm}|C{15mm}|C{70mm}|
}
\hline
\bfseries 系号 & \cell{\bfseries 学院} & \bfseries 系号 & \cell{\bfseries 学院/书院} \\
\hline
01 & \cell{材料科学与工程学院} & 25 & \cell{国际学院} \\
\hline
02 & \cell{电子信息工程学院} & 26 & \cell{新媒体艺术与设计学院} \\
\hline
03 & \cell{自动化科学与电气工程学院} & 27 & \cell{化学学院} \\
\hline
04 & \cell{能源与动力工程学院} & 28 & \cell{马克思主义学院} \\
\hline
05 & \cell{航空科学与工程学院} & 29 & \cell{人文与社会科学高等研究院} \\
\hline
06 & \cell{计算机学院} & 30 & \cell{空间与地球科学学院} \\
\hline
07 & \cell{机械工程及自动化学院} & 31 & \cell{无人系统研究院} \\
\hline
08 & \cell{经济管理学院} & 32 & \cell{航空发动机研究院} \\
\hline
09 & \cell{数学科学学院} & 33 & \cell{体育部} \\
\hline
10 & \cell{生物与医学工程学院} & 36 & \cell{国际前沿交叉科学研究院} \\
\hline
% 37 系北航学院分为七个书院
11 & \cell{人文社会科学学院（公共管理学院）} & 
\multirow{7}{*}{37} & \multirow{7}{*}{\shortstack{北\\航\\学\\院}} & 
71 & 传源书院（计算与智能科学类） \\
\cline{1-4}\cline{7-8}
12 & \cell{外国语学院} & & & 73 & 士谔书院（信息科学与技术类） \\
\cline{1-4}\cline{7-8}
13 & \cell{交通科学与工程学院} & & & 74 & 冯如书院（航空航天类） \\
\cline{1-4}\cline{7-8}
14 & \cell{可靠性与系统工程学院} & & & 75 & 士嘉书院（航空航天类） \\
\cline{1-4}\cline{7-8}
15 & \cell{宇航学院} & & & 76 & 守锷书院（航空航天类） \\
\cline{1-4}\cline{7-8}
16 & \cell{飞行学院} & & & 77 & 致真书院（理科试验班类） \\
\cline{1-4}\cline{7-8}
17 & \cell{仪器科学与光电工程学院} & & & 79 & 知行书院（社会科学试验班） \\
\hline
18 & \cell{北京学院} & 38 & \cell{医工交叉创新研究院} \\
\hline
19 & \cell{物理学院} & 39 & \cell{网络空间安全学院} \\
\hline
20 & \cell{法学院} & 41 & \cell{集成电路科学与工程学院} \\
\hline
21 & \cell{软件学院} & 42 & \cell{人工智能学院（人工智能研究院）} \\
\hline
22 & \cell{继续教育学院} & 45 & \cell{医学科学与工程学院} \\
\hline
% 23 系沈元学院分为四个学院
\multirow{4}{*}{23} & \multirow{4}{*}{\shortstack{沈\\元\\学\\院}} & 
231 & 高等理工学院 & 54 & \cell{大科学装置研究院/零磁科学中心} \\
\cline{3-8}
& & 232 & 未来空天技术学院 & 57 & \cell{杭州国际创新研究院} \\
\cline{3-8}
& & 233 & 国家卓越工程师学院 & 60 & \cell{中法航空学院} \\
\cline{3-8}
& & 234 & 量子科技学院 & 61 & \cell{中法未来科技学院} \\
\hline
24 & \cell{中法工程师学院/国际通用工程学院} & 62 & \cell{中西智能学院} \\
\hline
\end{tabular}

\end{center}
\end{document}